\documentclass[12pt, letterpaper]{article}
\author{Копылов Иван Александрович}
\usepackage[T1,T2A]{fontenc}
\usepackage[utf8]{inputenc}
\usepackage[english, russian]{babel}
\usepackage{graphicx}
\usepackage{booktabs}
\usepackage{listings}
\usepackage{a4, marginfit}
\graphicspath{{images/}}

\lstset{
    basicstyle=\fontencoding{T1}\ttfamily\small,
    breaklines=true,
    columns=fullflexible,
    keepspaces=true,
    frame=single
}

\title{Численные методы одномерной минимизации с использованием производной}
\date{Москва 2026}

\begin{document}
\maketitle
Лабораторная работа №2

Вариант 10

Копылов Иван Александрович, БИВТ-24-9\\

\includegraphics[width=2cm]{signature.jpg}

НИТУ МИСИС

Лычёв Андрей Владимирович

\newpage

\section{Цель работы}
Приобретение практических навыков для решения задач одномерной минимизации численными методами с использованием производной.

\section{Задача}
Требуется найти безусловный минимум функции одной переменной $y = f(x)$ на отрезке $[a, b]$, где функция является выпуклой и дифференцируемой. Необходимо найти такую точку $x^* \in [a, b]$, что
\[
    f(x^*) = \min\limits_{x \in [a, b]} f(x).
\]

Согласно варианту 10:
\[
    f(x) = 5x^2 + \frac{1}{x}, \quad
    a = 0.2, \quad
    b = 2, \quad
    \varepsilon = 0.0002.
\]

Для поиска минимума реализованы два метода:
\begin{itemize}
    \item \textit{Метод 1} --- метод средней точки;
    \item \textit{Метод 2} --- метод Брента.
\end{itemize}

\section{GenAI}
\textbf{Использованные инструменты GenAI:}
\begin{enumerate}
    \item OpenAI --- GPT-5. Использовался для оформления отчёта по данным из программного кода.
\end{enumerate}

\section{Графическое представление функции на заданном интервале}
\begin{center}
    \includegraphics[width=0.85\textwidth]{func.png}
\end{center}

\section{Аналитическое нахождение минимума}
Найдём производные функции:
\[
    f'(x) = \left(5x^2 + \frac{1}{x}\right)' = 10x - \frac{1}{x^2},
\]
\[
    f''(x) = 10 + \frac{2}{x^3}.
\]

Необходимое условие экстремума:
\[
    10x - \frac{1}{x^2} = 0,
    \quad
    10x^3 = 1,
    \quad
    x^* = \sqrt[3]{0.1} \approx 0.4641588834.
\]

Так как
\[
    f''(x^*) = 10 + \frac{2}{(x^*)^3} = 30 > 0,
\]
точка $x^*$ является точкой минимума.

Значение функции в точке минимума:
\[
    f(x^*) = 5(x^*)^2 + \frac{1}{x^*}
    \approx 3.2316520350.
\]

\section{Листинг программы}
\begin{lstlisting}[language=Python]
import math
import matplotlib.pyplot as plt
import numpy as np

def f_base(x):
    ans = 5 * x**2 + 1 / x
    return ans

def d_base(x):
    ans = 10*x - 1/x**2
    return ans

# graph
xs = np.linspace(0.2, 2, 100)
fs = np.array(list(map(f_base, xs)))

plt.plot(xs, fs)
plt.show()

# Method 1: midpoint method
def method_1(a, b, eps):
    iter_c = 0
    f_call_c = 0

    def f(x):
        nonlocal f_call_c
        f_call_c += 1
        return f_base(x)

    def d(x):
        nonlocal f_call_c
        f_call_c += 1
        return d_base(x)

    x = (a+b)/2
    d_x = d(x)

    while abs(a - b) > eps or abs(d_x) > eps:
        iter_c += 1
        if d_x > 0:
            b = x
        else:
            a = x
        x = (a+b)/2
        d_x = d(x)

    return (x, f(x), iter_c, f_call_c)

# Method 2: Brent method
def method_2(a, b, eps):
    iter_c = 0
    f_call_c = 0

    def f(x):
        nonlocal f_call_c
        f_call_c += 1
        return f_base(x)

    f_a = f(a)
    f_b = f(b)
    x = a + (3 - math.sqrt(5)) / 2 * (b - a)
    f_x = f(x)
    e = 0

    while max(x - a, b - x) > 2 * eps:
        iter_c += 1
        m = (a + b) / 2
        p = math.pow(b - x, 2) * (f_a - f_x) - math.pow(x - a, 2) * (f_b - f_x)
        q = 2 * ((b - x)*(f_a - f_x) + (x - a)*(f_b - f_x))
        s = x + p/q

        if (q == 0 and abs(e) < eps and (s < a or s > b)) or abs(p/q) > abs(e)/2:
            if x >= m:
                u = (math.sqrt(5) - 1) / 2 * x + (3 - math.sqrt(5)) / 2 * a
                f_u = f(u)
            else:
                u = (math.sqrt(5) - 1) / 2 * x + (3 - math.sqrt(5)) / 2 * b
                f_u = f(u)
        else:
            u = x + p/q
            f_u = f(u)
            if (abs(u - x) < 2*eps and u - a < 2*eps) or (b - u) < 2*eps:
                u = x + eps * np.sign(p/q)
                f_u = f(u)

        if f_u <= f_x:
            if u < x:
                b = x
            else:
                a = x
            x = u
            f_x = f_u
        else:
            if u < x:
                a = u
            else:
                b = u

    return (x, f_x, iter_c, f_call_c)

a = 0.2
b = 2
eps = 0.0002

(xm_1, f_xm_1, iter_c_1, f_call_c_1) = method_1(a, b, eps)
print(iter_c_1, f_call_c_1, xm_1, f_xm_1)

(xm_2, f_xm_2, iter_c_2, f_call_c_2) = method_2(a, b, eps)
print(iter_c_2, f_call_c_2, xm_2, f_xm_2)
\end{lstlisting}

\section{Результаты вычислений}
В программе были получены следующие результаты:
\begin{center}
   \small
   \resizebox{\textwidth}{!}{%
   \begin{tabular}{l c c c c}
        \toprule
        Название метода & Вычислений & Итераций & Результат вычислений & Погрешность по $x$ \\
        \midrule
        Аналитическое значение & $\sim$ & $\sim$ & $x^* \approx 0.4641588834, f(x^*) \approx 3.2316520350$ & $\sim$ \\
        Метод средней точки & 18 & 16 & $x = 0.4641525269$, $f(x) = 3.2316520357$ & $0.0000063565$ \\
        Метод Брента & 20 & 17 & $x = 0.4642476346$, $f(x) = 3.2316521532$ & $0.0000887513$ \\
        \bottomrule
    \end{tabular}%
   }
\end{center}


\section{Сравнительная характеристика методов}
Метод средней точки выполнил 16 итераций и 18 вычислений функции или производной. После исправления условия остановки метод дал точное приближение к аналитическому минимуму: погрешность по координате составила примерно $0.00000635$.

Метод Брента потребовал 17 итераций и 20 вычислений функции. Найденная точка $x = 0.4642476346$ также отличается от аналитического значения менее чем на $0.0001$, однако в данном эксперименте она оказалась немного менее точной по координате, чем результат метода средней точки.

Таким образом, в данном вычислительном эксперименте оба метода достигли заданной точности, но метод средней точки оказался немного эффективнее: он потребовал меньше итераций и вычислений, а также дал меньшую погрешность по $x$.

\section{Выводы}
В ходе лабораторной работы была исследована функция $f(x) = 5x^2 + 1/x$ на отрезке $[0.2, 2]$.
Аналитически установлено, что минимум достигается при
\[
    x^* \approx 0.4641588834,
    \quad
    f(x^*) \approx 3.2316520350.
\]

По данным программы оба метода дали результаты, близкие к аналитическому минимуму. Метод средней точки показал наилучший результат в этом эксперименте: при меньшем числе вычислений он получил более точное значение координаты минимума. Метод Брента также подтвердил сходимость к минимуму, но потребовал немного больше вычислений.

\end{document}
